\documentclass[12pt,a4paper]{article}
\usepackage{expediente}
\begin{document}
% =============================================================================
% FICHA DE DATOS DEL EXPEDIENTE --- actualizado 20 ago 2026
% =============================================================================

% --- VENDEDOR / TITULAR REGISTRAL -------------------------------------------
\newcommand{\VendedorNombres}{RICHAR DARWIN}
\newcommand{\VendedorApellidos}{CUTIRE CONDORI}
\newcommand{\VendedorNombreCompleto}{RICHAR DARWIN CUTIRE CONDORI}
\newcommand{\VendedorDNI}{42610690}
\newcommand{\VendedorEstadoCivil}{soltero, sin unión de hecho vigente}
\newcommand{\VendedorOcupacion}{\faltante{[ocupación de Richar]}}
\newcommand{\VendedorDomicilio}{\faltante{[domicilio actual de Richar]}}
\newcommand{\VendedorTelefono}{\faltante{[teléfono de Richar]}}
\newcommand{\VendedorCorreo}{\faltante{[correo]}}
\newcommand{\VendedorNacionalidad}{peruano}
\newcommand{\VendedorFechaNac}{\faltante{[fecha de nacimiento]}}
\newcommand{\VendedorDNIEmision}{\faltante{[fecha de emisión del DNI de Richar --- la pide SUNARP]}}

% --- CONVIVIENTE: NO CORRESPONDE -------------------------------------------
\newcommand{\ConvivienteNombre}{NO CORRESPONDE}
\newcommand{\ConvivienteDNI}{NO CORRESPONDE}
\newcommand{\ConvivienteDesde}{NO CORRESPONDE}
\newcommand{\ConvivienteInscrita}{no existe unión de hecho vigente}
\newcommand{\ConvivienteDomicilio}{NO CORRESPONDE}

% --- COMPRADORES (copropiedad en partes iguales) ---------------------------
\newcommand{\CompradorUnoNombre}{NICO ALVARO CUTIRE ARCE}
\newcommand{\CompradorUnoDNI}{48325017}
\newcommand{\CompradorUnoEmision}{01 de enero de 2024}
\newcommand{\CompradorUnoEstadoCivil}{soltero}
\newcommand{\CompradorUnoOcupacion}{ingeniero informático}
\newcommand{\CompradorUnoDomicilio}{Jr.\ Wiracocha, Sicuani, provincia de Canchis, Cusco, Perú}
\newcommand{\CompradorUnoTelefono}{984715108}

\newcommand{\CompradorDosNombre}{LUCY CUTIRE ARCE}
\newcommand{\CompradorDosDNI}{47478852}
\newcommand{\CompradorDosEmision}{09 de enero de 2026}
\newcommand{\CompradorDosEstadoCivil}{soltera}
\newcommand{\CompradorDosOcupacion}{contador}
\newcommand{\CompradorDosDomicilio}{Comunidad Hanocca, distrito de Layo, provincia de Canas, Cusco, Perú}
\newcommand{\CompradorDosTelefono}{953440875}

\newcommand{\CompradorNombres}{NICO ALVARO y LUCY}
\newcommand{\CompradorApellidos}{CUTIRE ARCE}
\newcommand{\CompradorNombreCompleto}{NICO ALVARO CUTIRE ARCE y LUCY CUTIRE ARCE}
\newcommand{\CompradorDNI}{48325017 y 47478852}
\newcommand{\CompradorEstadoCivil}{solteros}
\newcommand{\CompradorConyuge}{NO CORRESPONDE}
\newcommand{\CompradorConyugeDNI}{NO CORRESPONDE}
\newcommand{\CompradorOcupacion}{ingeniero informático y contador, respectivamente}
\newcommand{\CompradorDomicilio}{Sicuani (Canchis) y Layo (Canas), Cusco}
\newcommand{\CompradorTelefono}{984715108 / 953440875}
\newcommand{\CompradorCorreo}{\faltante{[correo de contacto]}}
\newcommand{\CompradorNacionalidad}{peruanos}
\newcommand{\PorcentajeCopropiedad}{50\,\% cada uno}

% --- INTERMEDIARIO / OCUPANTE DEL REMANENTE --------------------------------
\newcommand{\IntermediarioNombre}{LUCIO GIL CUTIRE PARI}
\newcommand{\IntermediarioDNI}{42467850}
\newcommand{\IntermediarioEstadoCivil}{\faltante{[estado civil de Lucio]}}
\newcommand{\IntermediarioOcupacion}{\faltante{[ocupación]}}
\newcommand{\IntermediarioDomicilio}{\faltante{[domicilio de Lucio]}}
\newcommand{\IntermediarioTelefono}{\faltante{[teléfono de Lucio]}}
\newcommand{\IntermediarioParentesco}{familiar (apellido Cutire)}

% --- PREDIO MATRIZ ----------------------------------------------------------
\newcommand{\PartidaMatriz}{\faltante{[N.º de partida electrónica SUNARP --- pendiente]}}
\newcommand{\OficinaRegistral}{Oficina Registral de Puerto Maldonado}
\newcommand{\AreaMatriz}{\faltante{[área total inscrita; las partes refieren del orden de 100 ha]}}
\newcommand{\UnidadCatastral}{\faltante{[código catastral / UC, si existe]}}
\newcommand{\NombreFundo}{\faltante{[nombre del fundo, si tiene]}}
\newcommand{\KmCarretera}{Interoceánica Sur, al este del Puente Jayave (el puente queda al oeste del predio, sobre el río; no debajo del lote)}

% --- PREDIO A INDEPENDIZAR --------------------------------------------------
\newcommand{\AreaIndependizarHa}{10{,}00}
\newcommand{\AreaIndependizarMetros}{100 000{,}00}
\newcommand{\PerimetroM}{2 200{,}58}
\newcommand{\FrenteM}{100{,}29}
\newcommand{\FondoM}{1 000{,}00}

\newcommand{\PUnoLat}{$-12{,}912452$}
\newcommand{\PUnoLon}{$-70{,}170058$}
\newcommand{\PDosLat}{$-12{,}912497$}
\newcommand{\PDosLon}{$-70{,}170981$}
\newcommand{\PTresLat}{$-12{,}903469$}
\newcommand{\PTresLon}{$-70{,}171438$}
\newcommand{\PCuatroLat}{$-12{,}903424$}
\newcommand{\PCuatroLon}{$-70{,}170515$}

\newcommand{\PUnoE}{373 065{,}00}
\newcommand{\PUnoN}{8 572 256{,}00}
\newcommand{\PDosE}{372 965{,}00}
\newcommand{\PDosN}{8 572 251{,}00}
\newcommand{\PTresE}{372 911{,}00}
\newcommand{\PTresN}{8 573 249{,}00}
\newcommand{\PCuatroE}{373 011{,}00}
\newcommand{\PCuatroN}{8 573 255{,}00}

\newcommand{\FechaPago}{15 de agosto de 2018}
\newcommand{\MontoPago}{S/\ 20 000{,}00}
\newcommand{\MontoPagoLetras}{VEINTE MIL Y 00/100 SOLES}
\newcommand{\FormaPagoHistorica}{dinero en efectivo, sin medio de pago del sistema financiero}

% Nombres genéricos --- sustituir por los reales al ir a notaría
\newcommand{\NotarioNombre}{Dr.\ JUAN CARLOS RÍOS SALINAS (nombre genérico a sustituir)}
\newcommand{\NotariaCiudad}{Puerto Maldonado}
\newcommand{\ColegiaturaAbogado}{CAL Madre de Dios N.º 0000 (genérico)}
\newcommand{\AbogadoNombre}{Abog.\ ANA MARÍA QUISPE HUAMÁN (nombre genérico a sustituir)}
\newcommand{\IngenieroNombre}{Ing.\ PEDRO ANTONIO MAMANI CONDORI (nombre genérico a sustituir)}
\newcommand{\IngenieroCIP}{CIP 000000 (genérico)}
\newcommand{\Municipalidad}{Municipalidad Distrital de Inambari}
\newcommand{\Provincia}{Tambopata}
\newcommand{\Departamento}{Madre de Dios}
\newcommand{\Distrito}{Inambari}

\newcommand{\LugarOtorgamiento}{Puerto Maldonado}
\newcommand{\DiaOtorgamiento}{\faltante{[día]}}
\newcommand{\MesOtorgamiento}{\faltante{[mes]}}
\newcommand{\AnioOtorgamiento}{2026}

\newcommand{\TestigoUno}{\faltante{[testigo 1: nombres, DNI, domicilio]}}
\newcommand{\TestigoDos}{\faltante{[testigo 2: nombres, DNI, domicilio]}}

\InicioDocumento{Modelo de cómo debió ser el documento privado de 2018}{DOC-22}

\noindent El papel de 2018 fue una hoja simple. Abajo van \textbf{dos textos}: (A) reconstrucción mínima para protocolizar lo que ocurrió; (B) el contrato que \textbf{debió} firmarse entonces. No reescriba el original antedatándolo: anexe foto del papel real y use el DOC-15 de reconocimiento.

\section*{A. Reconstrucción del trato (para anexar a la escritura de 2026)}
{\small
\begin{quote}
En Inambari / Mazuko, a los 15 días de agosto de 2018, \textbf{\VendedorNombreCompleto}, DNI \VendedorDNI, titular del predio rústico inscrito a su nombre, declara haber recibido de \textbf{\CompradorUnoNombre}\ y/o \textbf{\CompradorDosNombre}\ y/o por intermedio de \textbf{\IntermediarioNombre}, DNI \IntermediarioDNI, la suma de \textbf{\MontoPago} (\MontoPagoLetras) en \textbf{efectivo}, como precio de compraventa de \textbf{diez hectáreas} (10,00 ha) que se desprenderán del predio de mayor extensión ubicado en el distrito de Inambari, Tambopata, Madre de Dios, frente norte de la Interoceánica Sur, al este del Puente Jayave (el puente queda al oeste del lote, sobre el río).

El vendedor se obliga a independizar dichas diez hectáreas y a otorgar escritura pública a favor de los compradores. Lucio Gil Cutire Pari reconoce que el dinero de las 10 ha lo recibió el titular Richar Darwin Cutire Condori, y que Lucio no es dueño registral de ese recorte. El remanente del predio no es objeto de este pago.

Firmas: vendedor, compradores e intermediario. Huellas. DNI.
\end{quote}}

\section*{B. Cómo debió formalizarse en 2018 (contrato correcto)}
\begin{enumerate}[leftmargin=1.15cm]
\item Identificación completa de \textbf{vendedor} (titular SUNARP) y de \textbf{los dos compradores}, con DNI, estado civil y domicilio.
\item Descripción del inmueble: partida, área matriz, y las 10 ha con linderos o croquis. No basta decir ``diez hectáreas''.
\item Precio: S/\ 20\,000 y \textbf{depósito o transferencia} a la cuenta de Richar (Ley 28194). Conservar voucher.
\item Minuta autorizada por abogado y \textbf{escritura pública} el mismo año, o al menos firmas legalizadas ante notario.
\item Independización previa o simultánea, con plano UTM.
\item Alcabala y predial al día.
\item Inscripción en SUNARP a nombre de Nico y Lucy.
\item Si Lucio intermediaba, que firme como interviniente, no como vendedor, salvo que fuera el dueño (no lo era).
\end{enumerate}

\section*{C. Qué no hacer ahora}
No fabrique un contrato de 2018 con fecha vieja. No deposite S/\ 20\,000 en 2026 para ``bancarizar''. Use el reconocimiento DOC-15 + esta minuta 2026.

\vspace{0.4em}
\noindent\textit{Nota:} se pidió contrastar el documento adjunto; la imagen no llegó legible en este entorno. Cuando suba foto nítida del anverso y reverso, se transcribe cláusula por cláusula.

\LugarFecha
\CuatroFirmas{VENDEDOR}{\VendedorNombreCompleto}{COMPRADOR 1}{\CompradorUnoNombre}{COMPRADORA 2}{\CompradorDosNombre}{INTERMEDIARIO}{\IntermediarioNombre}

\end{document}
